\chapter{Código OpenQASM 3 de la demostración de Grover}
\label{ch:grover-qasm}

Este apéndice recoge el texto OpenQASM~3 completo que emite el compilador para \lstinline|examples/grover_2bit.qalc|, comentado y descrito en el Capítulo~\ref{ch:demostracion}. Se reproduce tal cual lo produce \lstinline|qalc-cli|, sin ninguna edición manual salvo el ajuste de los saltos de línea para que quepa en el ancho de la página.

\begin{lstlisting}[language={}, caption=Salida OpenQASM 3 completa de la demostración de Grover]
OPENQASM 3.0;
include "stdgates.inc";
qubit[2] q0;
qubit[2] q1;
qubit[2] q2;
qubit[2] q3;
qubit[1] q4;
qubit[1] q5;
qubit[1] q6;
output bit[2] c0;

h q0;
x q1[0];
x q2[0];
x q2[1];
cx q0[0], q3[0];
cx q1[0], q3[0];
ccx q0[0], q1[0], q4[0];
cx q0[1], q3[1];
cx q1[1], q3[1];
cx q4[0], q3[1];
ccx q0[1], q1[1], q5[0];
ccx q0[1], q4[0], q5[0];
ccx q1[1], q4[0], q5[0];
// uncompute q5
// uncompute q4
x q6;
h q6;
cx q3[0], q2[0];
x q2[0];
cx q3[1], q2[1];
x q2[1];
ccx q2[0], q2[1], q6[0];
x q2[1];
cx q3[1], q2[1];
x q2[0];
cx q3[0], q2[0];
// uncompute q6
// uncompute q2
// uncompute q3
// uncompute q1
h q0;
x q0;
cz q0[0], q0[1];
x q0;
h q0;
c0 = measure q0;
\end{lstlisting}
